\chapter{Writing Parallel Smalltalk}
\label{ch:parallel}

\standfirst{The classes 1983 has no equivalent of, the patterns that work, and
the mistakes that are specific to this system.}

\section{\texttt{forkParallel:} and \texttt{Promise}}

The one to reach for first. It forks a block and answers a \ct{Promise} for
its value:

\begin{st}
| p |
p := Processor forkParallel: [ (1 to: 1000000) inject: 0 into: [:a :b | a + b] ].
"...do something else..."
p value                          "waits if it is not ready yet"
\end{st}

The value comes back, and so does the failure: if the block raises, the
exception is re-signalled in whoever asks for the value---the process able to
do something about it.

\begin{center}
\small
\begin{tabular}{@{}ll@{}}
\toprule
\ctp{value} & the value, waiting if necessary; raises if it was broken\\
\ctp{valueOrNil} & the value, or nil if broken\\
\ctp{isSettled} \ctp{isBroken} & without waiting\\
\ctp{fulfill:} \ctp{signalError:} & from the producing side\\
\bottomrule
\end{tabular}
\end{center}

A promise settles once. A second fulfil or break is refused rather than
silently believed, because a value that changes after being read is not a
promise.

\subsection{Fan out over a collection}

\begin{st}
Processor parallelCollect: (1 to: 8) do: [:each | self expensiveWork: each]
\end{st}

Answers the results in order. Every promise is asked for, so a failure
anywhere arrives here rather than being lost. This is the whole of
map-parallel, and for a coarse-grained job it is all you need.

\section{\texttt{Mutex}}

\begin{st}
| lock |
lock := Mutex new.
lock critical: [ sharedThing add: x ]
\end{st}

Two things it does that a bare \ct{Semaphore} does not:

\begin{description}
\item[The release cannot be skipped.] It is in an \ct{ensure:}, so a non-local
  return or an exception through the critical block still unlocks. A bare
  semaphore leaves the lock held forever and every other worker waiting on
  nothing.
\item[Re-entry is detected.] A semaphore has no idea who holds it, so a
  process that takes the same one twice waits for itself, forever, with no
  diagnostic---the single most common way to deadlock a Smalltalk image. The
  owner is recorded, so that case raises an error naming it. It is not made to
  work: a recursive mutex is a different promise, and one that hides the design
  mistake rather than reporting it.
\end{description}

\begin{stnote}
There is deliberately no non-blocking \ct{tryCritical:}. It needs a wait that
can fail, and there is no primitive for one; writing it in Smalltalk means
testing \ct{excessSignals} and then decrementing---the exact test-then-act race
that was removed from \ct{Semaphore>>wait}.
\end{stnote}

\section{\texttt{Monitor}}

A reentrant lock with condition variables, for when "wait until something is
true" is the shape of the problem:

\begin{st}
| m |
m := Monitor new.
m critical: [ ... ].
m critical: [ m waitUntil: [ queue notEmpty ]. queue removeFirst ].
m critical: [ queue add: x. m signalAll ].
\end{st}

\begin{center}
\small
\begin{tabular}{@{}ll@{}}
\toprule
\ctp{critical:} & hold the monitor for the block; reentrant\\
\ctp{waitUntil:} & release, wait for a signal, re-acquire, re-test\\
\ctp{waitForChange} & release and wait for one signal\\
\ctp{signal} \ctp{signalAll} & wake one, or all\\
\ctp{hasWaiters} & \\
\bottomrule
\end{tabular}
\end{center}

\ct{waitUntil:} re-tests its condition after every wake, which is the correct
discipline and saves you from writing the loop wrong.

\section{\texttt{SharedQueue}}

Multi-producer, multi-consumer, blocking. This replaces the 1983
\ct{SharedQueue}, which guards an \ct{OrderedCollection} with two semaphores
--- correct when "concurrent" means green processes interleaved by one
interpreter, and wrong when eight native threads are inside it at once.

\begin{st}
| q |
q := SharedQueue new.
4 timesRepeat: [ [ [ q nextPut: self produce ] repeat ] fork ].
[ q next ]                        "blocks until something arrives"
\end{st}

\begin{center}
\small
\begin{tabular}{@{}ll@{}}
\toprule
\ctp{nextPut:} & add\\
\ctp{next} & remove, blocking\\
\ctp{nextOrNil} & remove, or nil if empty\\
\ctp{peek} \ctp{isEmpty} \ctp{size} & \\
\bottomrule
\end{tabular}
\end{center}

A queue is usually a better answer than a lock: it turns shared mutable state
into a handoff, which is \emph{reorganize} rather than \emph{serialize}.

\section{Patterns that work}

\subsection{Partition, compute, merge}

The one that scales. Give each worker its own data, and touch nothing shared
until the end:

\begin{st}
| parts results |
parts := self splitWorkInto: Processor workerCount.
results := Processor parallelCollect: parts do: [:part |
    part inject: 0 into: [:a :b | a + (self scoreOf: b)]].
^results inject: 0 into: [:a :b | a + b]
\end{st}

No lock anywhere, because nothing is shared. Each sub-block allocates from its
own worker's buffer.

\subsection{Producer--consumer}

A \ct{SharedQueue} between stages. Producers never see consumers' state and
the queue is the only shared object.

\subsection{Fork and join with failures}

\begin{st}
| promises |
promises := jobs collect: [:j | Processor forkParallel: [ self run: j ] ].
^promises collect: [:p | p value ]     "any failure arrives here"
\end{st}

\section{Mistakes specific to this system}

\begin{stwarn}[Lost update]
\begin{st}
count := count + 1.        "on a shared object: wrong"
Total := Total + 1.        "on a class variable: wrong"
\end{st}
Correct in every green-threaded Smalltalk, and a lost update here. Guard it,
or give each worker its own counter and sum at the end.
\end{stwarn}

\begin{stwarn}[Lazy initialisation]
\begin{st}
^Default ifNil: [Default := Foo new]     "two workers build two"
\end{st}
Forbidden. Initialise eagerly in a class-side \ct{initialize}.
\end{stwarn}

\begin{stwarn}[Sharing a collection]
\begin{st}
"two workers, one OrderedCollection"
results add: x                            "corrupts it"
\end{st}
The base collections are not synchronised. Use a \ct{SharedQueue}, or a
\ct{Mutex}, or one collection per worker. \ct{JSONObject}, \ct{JSONArray}
and \ct{DbSchemaGraph} are the three that guard themselves.
\end{stwarn}

\begin{stwarn}[Raising priority for atomicity]
It stops nobody on another core. It never worked here and it never will.
\end{stwarn}

\section{Proving it actually ran in parallel}

A test that merely does not crash proves very little. The discipline used
throughout this system's own suites is: \textbf{every worker computes something
only it can check}, so a fault arrives as a wrong answer rather than as a hope
that a sanitizer noticed.

\begin{st}
| results |
results := Processor parallelCollect: (1 to: 8) do: [:i |
    (1 to: 100000) inject: 0 into: [:a :b | a + (b * i)]].
"each element is a different, checkable number"
\end{st}

And to check that more than one worker was involved:

\begin{st}
(Processor parallelCollect: (1 to: 100) do: [:i | Processor activeWorkerIndex])
    asSet size                            "> 1 if it really spread"
\end{st}

\begin{stwarn}
That answers 1 under \ct{-run} and under \ct{-bootstrap -eval}, because
neither starts the worker pool---see Chapter~\ref{ch:contract}. It is the
parallel test suites and \ct{make bench} that run on many native threads
today. Write to the contract regardless: the code is the same either way, and
only the timing changes.
\end{stwarn}

\section{Measuring}

\ct{make bench} runs the scaling benchmark: four kernels, swept from one
worker to many, with the total work fixed.

\begin{center}
\small
\begin{tabular}{@{}lp{0.60\textwidth}@{}}
\toprule
arithmetic & the interpreter's loop, nothing allocated\\
mandelbrot & heavy arithmetic, almost no allocation --- the ceiling\\
intervals & pure interpretation: sends, blocks, a context per activation\\
collections & heavy allocation and collection --- the case a shared heap makes
  hardest\\
\bottomrule
\end{tabular}
\end{center}

If your own code does not scale, the question to ask is which of those four it
most resembles. Allocation-heavy code contends on the object table;
arithmetic-heavy code does not. \ct{doc/SCALING.md} has the analysis.

\begin{stnote}
Build with \ct{TSAN=1} and run your tests. ThreadSanitizer reports only races
that \emph{actually executed}, so it is necessary and not sufficient---but a
race it does report is real, and it will find them long before a customer
does.
\end{stnote}
